% GENERATED FILE. Edit canonical lessons, then run npm run generate:books.
% canonical-source-hash: fnv1a64:54541d874c2e3d76
% canonical-lessons: HI-C348-praman, HI-C348-sabitkarna, HI-C348-tathya, HI-C348-pukhta, HI-C348-arop

\chapter{Proof}
\label{ch:hi-proof}
\hlchaptermodality{348}

\begin{chapteropening}
\textbf{By the end of this chapter:} \emph{I can ask for proof and give it.}

Complete the last lesson of chapter 348: I can ask for proof and give it.
\end{chapteropening}

\section[\texorpdfstring{\hi{प्रमाण} (pramāṇ)}{प्रमाण (pramāṇ)}]{\hi{प्रमाण} (pramāṇ) — proof (formal), evidence}

\label{lesson:HI-C348-praman}



\begin{quote}
Before the new one: say the Hindi for to support, then the Hindi for to agree.
\end{quote}

\subsection*{You'll want to know: \hi{प्रमाण}}
\textbf{\hi{प्रमाण}} — \emph{pramāṇ} — \textquotedblleft{}proof (formal), evidence\textquotedblright{}.

\begin{quote}
\hi{तुम्हारे} \hi{पास} \hi{क्या} \hi{प्रमाण} \hi{है}?
\end{quote}

\begin{quote}
\emph{tumhāre pās kyā pramāṇ hai?}
\end{quote}

\begin{quote}
What evidence do you have?
\end{quote}

Masculine.

\begin{cousinweb}
\textbf{\hi{प्रमाण}} (\emph{pramāṇ}) is a tatsama, borrowed from Sanskrit \emph{pramāṇa}, \textquotedblleft{}a measure, a standard; proof\textquotedblright{} (Monier-Williams).
\end{cousinweb}

\begin{grammarlens}[title={What you've built}]
One more word for giving reasons and opinions.
\end{grammarlens}

\subsection*{Your turn}
\begin{itemize}
\raggedright
  \item \emph{Say it:} \emph{pramāṇ}
  \item \emph{Say it:} \emph{pramāṇ}, once more
  \item \emph{Say it:} \emph{tumhāre pās kyā pramāṇ hai?} — What evidence do you have?
  \item \emph{Recall:} say the Hindi for to support, then the Hindi for to agree, then say \emph{pramāṇ} again
\end{itemize}

\begin{tcolorbox}[breakable,colback=teal!4,colframe=teal!35!black,title={Before you move on}]
What does \hi{प्रमाण} mean? (\textquotedblleft{}Proof (formal), evidence\textquotedblright{}.) Say it once more.
\end{tcolorbox}
\section[\texorpdfstring{\hi{साबित} \hi{करना} (sābit karnā)}{साबित करना (sābit karnā)}]{\hi{साबित} \hi{करना} (sābit karnā) — to prove}

\label{lesson:HI-C348-sabitkarna}



\begin{quote}
Before the new one: say the Hindi for to agree, then the Hindi for proof (formal).
\end{quote}

\subsection*{You'll want to know: \hi{साबित} \hi{करना}}
\textbf{\hi{साबित} \hi{करना}} — \emph{sābit karnā} — \textquotedblleft{}to prove\textquotedblright{}.

\begin{quote}
\hi{मीरा} \hi{ने} \hi{अपनी} \hi{बात} \hi{साबित} \hi{की।}
\end{quote}

\begin{quote}
\emph{Mīrā ne apnī bāt sābit kī.}
\end{quote}

\begin{quote}
Mira proved her point.
\end{quote}

Past, with \hi{ने} (he or she …): \hi{उसने} \hi{अपनी} \hi{बात} \hi{साबित} \hi{की}, \emph{usne apnī bāt sābit kī}.

\begin{cousinweb}
\textbf{\hi{साबित} \hi{करना}} (\emph{sābit karnā}) joins \hi{साबित} to \hi{करना}, \textquotedblleft{}to do\textquotedblright{}. \hi{साबित} is usually traced to Arabic \emph{thābit}, \textquotedblleft{}firm, established\textquotedblright{}, which reached Hindi through Persian.
\end{cousinweb}

\begin{grammarlens}[title={What you've built}]
One more word for giving reasons and opinions.
\end{grammarlens}

\subsection*{Your turn}
\begin{itemize}
\raggedright
  \item \emph{Say it:} \emph{sābit karnā}
  \item \emph{Say it:} \emph{sābit karnā}, once more
  \item \emph{Say it:} \emph{Mīrā ne apnī bāt sābit kī.} — Mira proved her point.
  \item \emph{Recall:} say the Hindi for to agree, then the Hindi for proof (formal), then say \emph{sābit karnā} again
\end{itemize}

\begin{tcolorbox}[breakable,colback=teal!4,colframe=teal!35!black,title={Before you move on}]
What does \hi{साबित} \hi{करना} mean? (\textquotedblleft{}To prove\textquotedblright{}.) Say it once more.
\end{tcolorbox}
\section[\texorpdfstring{\hi{तथ्य} (tathya)}{तथ्य (tathya)}]{\hi{तथ्य} (tathya) — a fact}

\label{lesson:HI-C348-tathya}



\begin{quote}
Before the new one: say the Hindi for proof (formal), then the Hindi for to prove.
\end{quote}

\subsection*{You'll want to know: \hi{तथ्य}}
\textbf{\hi{तथ्य}} — \emph{tathya} — \textquotedblleft{}a fact\textquotedblright{}.

\begin{quote}
\hi{यह} \hi{राय} \hi{नहीं}, \hi{तथ्य} \hi{है।}
\end{quote}

\begin{quote}
\emph{yah rāy nahīṃ, tathya hai.}
\end{quote}

\begin{quote}
This is not an opinion; it is a fact.
\end{quote}

Masculine.

\begin{cousinweb}
\textbf{\hi{तथ्य}} (\emph{tathya}) is a tatsama, borrowed from Sanskrit \emph{tathya}, \textquotedblleft{}true; the truth\textquotedblright{} (Monier-Williams).
\end{cousinweb}

\begin{grammarlens}[title={What you've built}]
One more word for giving reasons and opinions.
\end{grammarlens}

\subsection*{Your turn}
\begin{itemize}
\raggedright
  \item \emph{Say it:} \emph{tathya}
  \item \emph{Say it:} \emph{tathya}, once more
  \item \emph{Say it:} \emph{yah rāy nahīṃ, tathya hai.} — This is not an opinion; it is a fact.
  \item \emph{Recall:} say the Hindi for proof (formal), then the Hindi for to prove, then say \emph{tathya} again
\end{itemize}

\begin{tcolorbox}[breakable,colback=teal!4,colframe=teal!35!black,title={Before you move on}]
What does \hi{तथ्य} mean? (\textquotedblleft{}A fact\textquotedblright{}.) Say it once more.
\end{tcolorbox}
\section[\texorpdfstring{\hi{पुख़्ता} (pukhtā)}{पुख़्ता (pukhtā)}]{\hi{पुख़्ता} (pukhtā) — firm, conclusive}

\label{lesson:HI-C348-pukhta}



\begin{quote}
Before the new one: say the Hindi for to prove, then the Hindi for a fact.
\end{quote}

\subsection*{You'll want to know: \hi{पुख़्ता}}
\textbf{\hi{पुख़्ता}} — \emph{pukhtā} — \textquotedblleft{}firm, conclusive\textquotedblright{}.

\begin{quote}
\hi{पुलिस} \hi{को} \hi{पुख़्ता} \hi{प्रमाण} \hi{चाहिए।}
\end{quote}

\begin{quote}
\emph{pulis ko pukhtā pramāṇ cāhiye.}
\end{quote}

\begin{quote}
The police need conclusive proof.
\end{quote}

\begin{cousinweb}
\textbf{\hi{पुख़्ता}} (\emph{pukhtā}) is usually traced to Persian \emph{pukhta}, \textquotedblleft{}cooked, ripe; firm\textquotedblright{}.
\end{cousinweb}

\begin{grammarlens}[title={What you've built}]
One more word for giving reasons and opinions.
\end{grammarlens}

\subsection*{Your turn}
\begin{itemize}
\raggedright
  \item \emph{Say it:} \emph{pukhtā}
  \item \emph{Say it:} \emph{pukhtā}, once more
  \item \emph{Say it:} \emph{pulis ko pukhtā pramāṇ cāhiye.} — The police need conclusive proof.
  \item \emph{Recall:} say the Hindi for to prove, then the Hindi for a fact, then say \emph{pukhtā} again
\end{itemize}

\begin{tcolorbox}[breakable,colback=teal!4,colframe=teal!35!black,title={Before you move on}]
What does \hi{पुख़्ता} mean? (\textquotedblleft{}Firm, conclusive\textquotedblright{}.) Say it once more.
\end{tcolorbox}
\section[\texorpdfstring{\hi{आरोप} (ārop)}{आरोप (ārop)}]{\hi{आरोप} (ārop) — an accusation, a charge}

\label{lesson:HI-C348-arop}



\begin{quote}
Before the new one: say the Hindi for a fact, then the Hindi for conclusive.
\end{quote}

\subsection*{You'll want to know: \hi{आरोप}}
\textbf{\hi{आरोप}} — \emph{ārop} — \textquotedblleft{}an accusation, a charge\textquotedblright{}.

\begin{quote}
\hi{पुलिस} \hi{का} \hi{आरोप} \hi{है} \hi{कि} \hi{ड्राइवर} \hi{ने} \hi{नियम} \hi{तोड़ा।}
\end{quote}

\begin{quote}
\emph{pulis kā ārop hai ki ḍrāivar ne niyam toṛā.}
\end{quote}

\begin{quote}
The police allege that the driver broke the rules.
\end{quote}

Masculine.

\begin{cousinweb}
\textbf{\hi{आरोप}} (\emph{ārop}) is a tatsama, borrowed from Sanskrit \emph{āropa}, \textquotedblleft{}imposing a burden, charging with\textquotedblright{} (Monier-Williams).
\end{cousinweb}

\begin{grammarlens}[title={What you've built}]
One more word for giving reasons and opinions.
\end{grammarlens}

\subsection*{Your turn}
\begin{itemize}
\raggedright
  \item \emph{Say it:} \emph{ārop}
  \item \emph{Say it:} \emph{ārop}, once more
  \item \emph{Say it:} \emph{pulis kā ārop hai ki ḍrāivar ne niyam toṛā.} — The police allege that the driver broke the rules.
  \item \emph{Recall:} say the Hindi for a fact, then the Hindi for conclusive, then say \emph{ārop} again
\end{itemize}

\begin{tcolorbox}[breakable,colback=teal!4,colframe=teal!35!black,title={Before you move on}]
What does \hi{आरोप} mean? (\textquotedblleft{}An accusation, a charge\textquotedblright{}.) Say it once more.
\end{tcolorbox}
